% llm_peer/drafts/zh/manuscript.tex
% P4.2 中文 manuscript 骨架 (plan-r3 衍生论文扩展, 2026-05-15)
% 编译: xelatex + biber + xelatex × 2

% papers_assets/_shared/preamble_zh.tex
% Shared Chinese LaTeX preamble (XeLaTeX + CJKutf8).
% P4.1 (plan-r3 derivative-paper extension, 2026-05-15)
% Compile: xelatex + biber + xelatex × 2

\documentclass[12pt,a4paper]{article}

% --- CJK Fonts (XeLaTeX, cross-platform Q1.1b fallback) ---
\usepackage{fontspec}
\usepackage{xeCJK}
% Windows (local dev): SimSun/SimHei/FangSong
% Linux/macOS (CI): Noto Serif CJK SC / Noto Sans CJK SC (fonts-noto-cjk apt package)
% Last resort: AR PL UMing CN (fonts-arphic-uming) which is preinstalled in many texlive distros
\IfFontExistsTF{SimSun}{%
  \setCJKmainfont{SimSun}[BoldFont=SimHei]%
}{%
  \IfFontExistsTF{Noto Serif CJK SC}{%
    \setCJKmainfont{Noto Serif CJK SC}[BoldFont=Noto Serif CJK SC]%
  }{%
    \setCJKmainfont{AR PL UMing CN}[AutoFakeBold=2]%
  }%
}
\IfFontExistsTF{SimHei}{%
  \setCJKsansfont{SimHei}%
}{%
  \IfFontExistsTF{Noto Sans CJK SC}{%
    \setCJKsansfont{Noto Sans CJK SC}%
  }{%
    \setCJKsansfont{AR PL UKai CN}%
  }%
}
\IfFontExistsTF{FangSong}{%
  \setCJKmonofont{FangSong}%
}{%
  \IfFontExistsTF{Noto Sans Mono CJK SC}{%
    \setCJKmonofont{Noto Sans Mono CJK SC}%
  }{%
    \setCJKmonofont{AR PL UMing CN}%
  }%
}
\setmainfont{Times New Roman}[Ligatures=TeX]
\xeCJKsetup{CJKecglue={\hskip 0.15em plus 0.05em minus 0.05em}}

% --- Math ---
\usepackage{amsmath, amssymb, amsthm}

% --- Tables & Figures ---
\usepackage{booktabs}
\usepackage{graphicx}
\usepackage{longtable}
\usepackage{rotating}
\usepackage{caption}
\captionsetup{font=small,labelfont=bf,skip=4pt}

% --- Layout ---
\usepackage[margin=1in]{geometry}
\usepackage{setspace}
\onehalfspacing
\usepackage{hyperref}
\hypersetup{colorlinks=true,linkcolor=blue,citecolor=blue,urlcolor=blue}

% --- Bibliography ---
\usepackage[backend=biber,style=apa,natbib=true]{biblatex}

% --- Macros ---
\newcommand{\code}[1]{\texttt{#1}}
\usepackage[dvipsnames]{xcolor}
\newcommand{\NoteThis}[1]{\textbf{\textcolor{red}{[TODO: #1]}}}

% --- Theorem environments ---
\newtheorem{prop}{命题}
\newtheorem{hyp}{假设}


\addbibresource{refs.bib}

\title{LLM 生成的同业识别与银行金融科技战略的因果推断}
\author{\NoteThis{作者姓名}}
\date{\today}

\begin{document}
\maketitle

\begin{abstract}
\NoteThis{最终摘要待定。}
本文基于银行—年度金融科技媒体报道, 采用LLM 定义的同业集 (Cao \emph{et~al.}\ 2025) + 留一法工具变量,
度量中国上市银行的LLM 同业留一法语调 IV (\code{llm\_peer\_tone\_iv\_full})。我们发现 \NoteThis{核心结论}\@。
\end{abstract}

\section{引言}\label{sec:intro}
\NoteThis{阐释研究问题与贡献。}

\section{文献综述}\label{sec:lit}
(1)~基于 LLM 的同业识别 (Cao \emph{et~al.}\ 2025 \emph{RAS}); (2)~Stock--Yogo (2005) 弱工具检验; (3)~Kleibergen--Paap (2006) 秩识别。

\section{方法}\label{sec:method}
\subsection{样本与数据}
银行—年度面板覆盖约 164 个银行—年观测 (2017--2018)。
所有数据由 \code{news\_process} 管线 (Phase 9b 基线) 处理,
详见 \code{papers\_assets/V6\_METHODOLOGY\_CHANGES.md}。

\subsection{语调标注 (F1: 三家族 LLM 真三角验证)}
采用三个独立 LLM 家族: \code{gpt-5.4-nano} (OpenAI)、
\code{gemini-3.1-flash-lite-preview} (Google)、\code{deepseek-v4-flash}
(DeepSeek\,/\,Grok)。统一通过 yunwu.ai (\code{/v1/chat/completions},
OpenAI 兼容) 调用, \code{temperature=0}, \code{seed=42}。
家族间一致性以 Cohen's $\kappa$ 与 Krippendorff's $\alpha$ 报告;
多数投票得最终标签, 平局标 \code{undefined} (非 \code{neutral})。

\subsection{篇章级编码 (F2: 去污染)}
三家族全部失败 / 三方分裂者编为\emph{缺失} (非 \code{neutral});
银行--金融科技相关性未通过多数验证者亦视为缺失。
此举消除了 v6 之前 81.7\,\% ``中性'' 标签的污染。

\subsection{零观测语义 (F3: NaN 主路径)}
对 \code{n\_articles=0} 的银行—年, 相似度类特征
(\code{auth\_sim\_*}、\code{benefit\_*}、\code{article\_sim\_*})、
比例类 (\code{token\_adopt\_*}、\code{hidden\_adopt\_*}、
\code{subst\_adopt\_*})、\code{fintech\_intensity\_mean} 均编为
\code{NaN}; 计数类保持 0。回归采用列删除 (Allison 2001 / Heckman 1979)。

\subsection{Bartik IV 修剪 (F4)}
约 30 个候选 Bartik 工具变量中, 28 个在 v6 中因违反
Goldsmith-Pinkham \emph{et~al.}\ (2020, \emph{AER}) 外生性被删除。
保留: \code{bartik\_type\_iv\_full} 与 \code{prov\_leave1out\_tone\_full}。

\section{数据描述}\label{sec:data}
\NoteThis{Table 1 描述统计参见 \code{papers/llm_peer/tables.py:build\_table1}。}

\section{基准回归}\label{sec:main}
\NoteThis{Table 2 源自 \code{output/tables/table2\_*.csv}; Figure 2 系数图。}

\section{机制检验}\label{sec:mech}
\NoteThis{Table 3 源自 \code{output/tables/table3\_*.csv}; Figure 3。}

\section{异质性}\label{sec:het}
\NoteThis{Table 4 (SOE × 论文专属交互) 源自
\code{output/tables/table4\_heterogeneity.csv}。}

\section{稳健性}\label{sec:robust}
\NoteThis{Table 5 (论文专属稳健性矩阵) 源自
\code{output/tables/table5\_*.csv}。}

\section{结论}\label{sec:concl}
\NoteThis{最终讨论与政策含义。}

\printbibliography

\end{document}
